\section{Investment Problem} \label{sec:investment_problem} \label{sec:investment-problem}

We consider the canonical dynamic investment problem of a single firm (firm $j$) that
serves as the standard model of firm-level capital allocation in macroeconomics.
The firm operates a production function $f(z_{j,t}, k_{j,t})$, concave and increasing
in both arguments, where $z_{j,t}$ is an idiosyncratic productivity shock governing
the firm's current efficiency and $k_{j,t}$ is its capital stock, the physical
machinery and equipment at its disposal.

The firm enters the period with a stock of debt $b_{j,t}$ which carries a risk-free
interest rate $r$. Each period, the firm earns a dividend:
\begin{equation}
    d_{j,t} = f(z_{j,t}, k_{j,t}) - i_{j,t} - \psi(i_{j,t}, k_{j,t})
    + b_{j,t+1} - (1+r)b_{j,t}.
\end{equation}
This represents the residual of output after investment expenditure $i_{j,t}$,
capital adjustment costs $\psi(i_{j,t}, k_{j,t})$, and the net cash flow from
issuing new debt $b_{j,t+1}$ and repaying existing debt obligations $(1+r)b_{j,t}$.
The firm derives a per-period payoff $u(d_{j,t})$ from this dividend. The firm's
problem is to choose policies for investment and debt that maximize the expected
present discounted value of its payoff stream. It is common to model firms as
risk-neutral agents with a linear payoff function, $u(d_{j,t}) = d_{j,t}$.

\begin{equation}
    \max_{\{i_{j,t},\, b_{j,t+1}\}_{t=0}^{\infty}} \;
    \mathbb{E}_0 \left[ \sum_{t=0}^{\infty}
    \beta^{t} d_{j,t} \right]
\end{equation}

The firm faces several frictions. First, a capital adjustment cost
$\psi(i_{j,t}, k_{j,t})$ penalizes rapid changes in the capital stock;
installing or dismantling capital is costly and disruptive. We consider a
standard quadratic adjustment cost,
\begin{equation}
    \psi(i_{j,t}, k_{j,t}) =
    \frac{\phi}{2} \left(\frac{i_{j,t}}{k_{j,t}}\right)^2 k_{j,t},
\end{equation}
where $\phi \geq 0$. This cost is convex in the rate of investment; the firm
cannot freely jump to its desired capital stock and must smoothly adjust its
capital stock. Second, a collateral borrowing constraint limits the firm's
ability to finance investment externally. The firm can only issue new debt
$b_{j,t+1}$ up to a fraction $\theta$ of its next-period capital stock:
\begin{equation}
    b_{j,t+1} \leq \theta k_{j,t+1}.
\end{equation}
This bounds the feasible investment action at each state and introduces the
possibility that profitable investment opportunities go unexploited due to
insufficient internal funds.

The capital stock depreciates over time at a per-period rate $\delta$, evolving
according to a standard law of motion:
\begin{equation}
    k_{j,t+1} = (1 - \delta) k_{j,t} + i_{j,t}.
\end{equation}

The firm's problem is therefore a dynamic programming problem with continuous
state and action spaces. Defining the state as the firm's idiosyncratic
productivity, capital, and debt $(z_{j,t}, k_{j,t}, b_{j,t})$, the problem can
be expressed as the recursive Bellman equation:
\begin{equation}
    V(z_{j,t}, k_{j,t}, b_{j,t}) =
    \max_{i_{j,t},\, b_{j,t+1}} \left\{
    d_{j,t} + \beta \, \mathbb{E}\!\left[
    V(z_{j,t+1}, k_{j,t+1}, b_{j,t+1})
    \,\middle|\, z_{j,t} \right] \right\}
\end{equation}
subject to:
\begin{align}
    d_{j,t} &= f(z_{j,t}, k_{j,t}) - i_{j,t}
    - \frac{\phi}{2} \left(\frac{i_{j,t}}{k_{j,t}}\right)^2 k_{j,t}
    + b_{j,t+1} - (1+r)b_{j,t} \\
    k_{j,t+1} &= (1 - \delta) k_{j,t} + i_{j,t} \\
    b_{j,t+1} &\leq \theta k_{j,t+1} \\
    d_{j,t} &\geq 0
\end{align}
where $V(\cdot)$ is the value function and the expectation is taken over the
stochastic evolution of productivity conditional on the current state.

\subsection{Solution to Canonical Investment Problem} \label{subsec:canonical-investment-solution}

The firm's optimal investment policy can be summarized using the shadow value of
installed capital ($q_{j,t}$). When the collateral constraint is slack,
the optimal investment-capital ratio is linear in $(q_{j,t}-1)$, where $q_{j,t}$
denotes the marginal value of installed capital. When the collateral constraint
binds, however, investment is limited by the firm's borrowing capacity and internal
funds. The resulting optimal investment policy can therefore be written as the following.
\begin{equation}
\frac{i_{j,t}}{k_{j,t}} =
\begin{cases}
\frac{1}{\phi}(q_{j,t}-1),
& \text{if } b_{j,t+1} < \theta k_{j,t+1} \\
\frac{k_{j,t+1}^{BC} - (1-\delta)k_{j,t}}{k_{j,t}},
& \text{if } b_{j,t+1} = \theta k_{j,t+1}
\end{cases}
\end{equation}
The term $k_{j,t+1}^{BC}$ denotes the level of next-period capital implied by the
binding borrowing constraint and the firm's budget constraint. Intuitively, when
the value of an additional unit of capital exceeds its installation cost, the firm
increases investment. When the collateral
constraint binds, capital also has value because it relaxes the borrowing limit,
and investment becomes limited by available collateral rather than the marginal
$q$ condition. The full derivation of the firm's first-order conditions and the
resulting Tobin's $q$ investment rule is provided in \ref{app:canonical-investment-derivation}.

\subsection{The Firm's Problem and Policy Under Specific Functional Forms} \label{subsec:canonical-investment-specific-forms}

To this point, we have specified the firm's investment problem and derived the firm's optimal investment policy in general terms.
For simplicity, we choose production to be Cobb-Douglas, $f(z_{j,t}, k_{j,t}) = z_{j,t} k_{j,t}^\alpha$, and the productivity shock to follow an AR(1) 
process, $\ln z_{j,t+1} = \rho \ln z_{j,t} + \varepsilon_{j,t+1}$ with $\varepsilon_{j,t+1} \sim \mathcal{N}(0, \sigma_z^2)$. 
We also assume linear utility, $u(d_{j,t}) = d_{j,t}$, so that the firm is risk-neutral.
Under these assumptions, the Bellman equation characterizing the firm's problem takes the following explicit form, 
subject to the capital accumulation law of motion and the collateral constraint, as described previously.
\begin{equation}
\begin{split}
    V(z_{j,t}, k_{j,t}, b_{j,t}) &= 
    \max_{i_{j,t},\, b_{j,t+1}} \Bigg\{ 
    z_{j,t} k_{j,t}^\alpha - i_{j,t} 
    - \frac{\phi}{2} \left(\frac{i_{j,t}}{k_{j,t}}\right)^2 k_{j,t} \\
    &\quad + b_{j,t+1} - (1+r)b_{j,t} 
    + \beta \, \mathbb{E} \left[ 
    V(z_{j,t+1}, k_{j,t+1}, b_{j,t+1}) 
    \,\middle|\, z_{j,t} \right] \Bigg\}
\end{split}
\end{equation}

Intuitively, the value function represents the fundamental value of the firm given its 
current state (defined by its productivity, capital, and debt $(z_{j,t}, k_{j,t}, b_{j,t})$) 
under a specific set of policies. It decomposes this value into two parts: the current 
dividend and the discounted expected value of future payoffs. The current dividend is 
simply the firm's production revenue, $z_{j,t} k_{j,t}^\alpha$, minus the total cost of 
investment including adjustment frictions, $i_{j,t} + \frac{\phi}{2} \left(\frac{i_{j,t}}{k_{j,t}}\right)^2 k_{j,t}$, 
plus the net cash flow from debt management, $b_{j,t+1} - (1+r)b_{j,t}$. Consequently, the 
maximized value function, $V(\cdot)$, corresponds to the fundamental equity value of 
the firm when it optimally executes its investment and financing decisions. Because the 
dividend represents the residual cash flow to owners after all operational and debt obligations are settled, maximizing this value function is 
economically equivalent to maximizing shareholder value.

Taking the first-order condition with respect to investment, the optimal policy relies 
on the shadow value of installed capital, $q_{j,t}$, which is explicitly defined as the 
expected marginal continuation value of an additional unit of capital:
\begin{equation}
    q_{j,t} = \beta \, \mathbb{E} \left[ \frac{\partial V(z_{j,t+1}, k_{j,t+1}, b_{j,t+1})}{\partial k_{j,t+1}} \,\middle|\, z_{j,t} \right].
\end{equation}
Using this definition, the firm's optimal investment-to-capital ratio can be written as:
\begin{equation}
\frac{i_{j,t}}{k_{j,t}} =
\begin{cases}
\frac{1}{\phi}(q_{j,t}-1),
& \text{if } b_{j,t+1} < \theta k_{j,t+1} \\[10pt]
\frac{k_{j,t+1}^{BC} - (1-\delta)k_{j,t}}{k_{j,t}},
& \text{if } b_{j,t+1} = \theta k_{j,t+1}
\end{cases}
\end{equation}

To visualize the rational firm's optimal policy, we show canonical phase diagrams
in \ref{fig:rational_firm_policy_phase_diagrams}: the optimal next-period capital
stock $k_{t+1}$ as a function of current capital $k_t$ for different productivity
levels $z_t$, and the corresponding marginal value of capital
$q_t$\footnote{Equivalently, the shadow price of capital or the Lagrange multiplier
on the capital accumulation constraint.} as a function of $k_t$. Together these
characterize the firm's policy: when $q_t > 1$ the firm invests and grows its
capital stock; when $q_t < 1$ it disinvests. The firm invests such that
$k_{t+1} = i_t + (1-\delta)k_t$ and $q_t = 1$ when the collateral constraint is
slack, and up to the borrowing limit when it binds. Because the rational firm knows
the optimal policy under any state $(z_t, k_t, b_t)$, it almost always avoids a
binding collateral constraint. A large negative productivity shock followed by a
large positive shock can produce a binding constraint, but this occurs with low 
probability under the ergodic distribution.
\begin{figure}[ht]
    \centering
    \includegraphics[width=1.0\textwidth]{../figures/rational_firm_policy_phase_diagrams.pdf}
    \caption{Phase diagrams of the rational firm's optimal policy.}
    \label{fig:rational_firm_policy_phase_diagrams}
\end{figure}

Notice in \ref{fig:rational_firm_policy_phase_diagrams} that the $q = 1$ crossings 
in panel (b) do not align with the steady-state crossings in panel (a). This is 
because $q = 1$ marks the extensive margin of investment, the capital stock at 
which the firm is indifferent between adding and not adding a unit of capital, 
i.e.\ $i = 0$. The $45°$ crossing marks the intensive margin, the capital stock 
at which the firm invests exactly enough to offset depreciation, $i = \delta k$. 
At the $q = 1$ crossing, the firm switches from divestment ($i < 0$, allowing 
capital to decay) to active investment ($i > 0$). At the $45°$ crossing, the firm 
reaches its stationary capital stock where gross investment equals depreciation. 
The two are separated by the adjustment cost wedge: at steady state, 
$q^* = 1 + \kappa \delta > 1$, so the firm must value capital strictly above 
replacement cost to justify the flow cost of maintaining it.

\subsection{Investment Problem Under Learning Through Reasoning and Experience}

With the canonical formulation in hand, we can now transition to the RL framework.
We now write the firm's problem using the reasoning and experience framework described
in \ref{app:experience_reasoning_agent_consumption}. Importantly, the firm exists
in the environment described previously (\ref{sec:investment-problem}), but 
it does not know the solution (\ref{subsec:canonical-investment-solution}).
The firm attempts to learn about its optimal policy through two processes: experience
and planning. 

The experience channel is a temporal-difference learning process with a Bayesian 
update governed by a Kalman gain, formally described in 
\ref{def:experience_learning_update}. Because the environment is stochastic, 
the firm may misattribute outcomes to actions. For example, a particularly adverse productivity 
realization can cause the firm to infer that its investment was too aggressive, 
while a favorable shock can reinforce excessive divestment. The reasoning channel 
resembles a model-based reinforcement learning process. At a cognitive cost, the 
firm mentally simulates candidate actions and their consequences, refining its 
beliefs about the action-value function beyond what experience alone provides. 
This is formally described in \ref{def:reasoning_learning_update}.

The two channels are linked through a certainty equivalence assumption: the firm 
chooses its reasoning intensity under the expected realization of the reasoning 
signals, so that reasoning reduces the posterior variance of $Q^*$ but does not 
shift the conditional mean. The cognitive cost of reasoning is proportional to 
the entropy reduction achieved, weighted by the shadow price of the entropy 
constraint (the Lagrange multiplier on the exploration bound). This shadow price 
captures the opportunity cost of uncertainty: when forced exploration is costly, 
variance reduction is worth the cognitive expenditure. Because the reasoning step 
preserves the conditional mean by construction, the belief sequence 
$\{\hat{Q}_t\}$ is a martingale with respect to the reasoning filtration. 
The firm cannot expect its own future reasoning to correct current errors, so 
deviations from the rational policy are persistent.

At each time $t$, the firm exists in state defined by its capital stock $k_t$, debt $b_t$, 
and productivity $z_t$ ($\mathbf{s}_t = (z_t, k_t, b_t)$). The firm chooses both investment
action $i_t$ and borrowing action $b_{t+1}$ ($\mathbf{a}_t = (i_t, b_{t+1})$) which 
determine the next-period state through the transition function $\mathcal{F}(i_t, b_{t+1}, (z_t, k_t, b_t), \varepsilon_{t+1})$.
\begin{equation}
\mathbf{s}_{t+1} = 
\begin{bmatrix}
        z_{t+1} \\
        k_{t+1} \\
        b_{t+1}
\end{bmatrix} = \mathcal{F}\left((i_t, b_{t+1}), (z_t, k_t, b_t), \varepsilon_{t+1}\right) =
\begin{bmatrix}
\exp\left(\rho \ln z_t + \varepsilon_{t+1}\right) \\
(1-\delta)k_t + i_t \\
b_{t+1}
\end{bmatrix}
\end{equation}

The optimal policy is retrieved by solving the firm's constrained optimization problem 
conditional on the state $\mathbf{s}_t = (z_t, k_t, b_t)$. In this context, 
the \textit{Value Function} ($V$) represents the maximized fundamental equity value of 
the firm, and the \textit{Action-Value Function} ($Q$) characterizes the immediate 
dividend from a specific investment and financing choice plus the expected discounted 
future value.

\begin{proposition}[Optimal Policy and Value Functions]
    The optimal policy $\pi^*$ maps the state $\mathbf{s}_t = (z_t, k_t, b_t)$ to actions $\mathbf{a}_t = (i_t, b_{t+1})$ subject to the collateral borrowing constraint. The Action-Value Function $Q^*$ and Value Function $V^*$ satisfy:
\begin{align}
        Q^*(\mathbf{s}_t, \mathbf{a}_t) &= z_t k_t^\alpha - i_t - \frac{\phi}{2} \left(\frac{i_t}{k_t}\right)^2 k_t + b_{t+1} - (1+r)b_t + \beta \mathbb{E} \left[ V^*(z_{t+1}, k_{t+1}, b_{t+1}) \,\middle|\, z_t \right] \\
        V^*(\mathbf{s}_t) &= \max_{\mathbf{a}_t \in \Gamma(\mathbf{s}_t)} Q^*(\mathbf{s}_t, \mathbf{a}_t) = Q^*(\mathbf{s}_t, \pi^*(\mathbf{s}_t))
\end{align}
where $\Gamma(\mathbf{s}_t)$ represents the feasible action space defined by the collateral constraint $b_{t+1} \leq \theta ((1-\delta)k_t + i_t)$ and the non-negative dividend condition $d_t \geq 0$.
\end{proposition}

The firm is Bayesian and has beliefs over the true action-value function $Q^*$ described by
a Gaussian Process with mean $\hat{Q}_t$ and covariance $\hat{\Sigma}_t$, identical to
\ref{subsec:agent_beliefs}. As described in \ref{subsec:information_directed_policies}, 
the firm's \textit{true} problem is a Bayes-adaptive Markov decision 
process (BAMDP) where the true $Q^*$ is unknown and beliefs continuously evolve. 
We assume the firm cannot solve it perfectly.\footnote{This problem
is also computationally intractable and cannot be solved in this continuous state and action space setting.}
Instead, the firm approximates the optimal policy over time through imperfect reasoning and experience.
The experience and reasoning cycle each period as described in \ref{def:learning_cycle}.

\begin{proposition}[Investment as the Sole Dynamic Choice Variable]
\label{prop:debt_separation}
The firm's state can be reduced from $\mathbf{s}_t = (z_t, k_t, b_t)$ to $\mathbf{s}_t = (z_t, k_t)$ by 
substituting two case optimal borrowing policy. The reduced Bellman equation is:
\begin{equation}
\begin{split}
    V^*(z_t, k_t) = \max_{i_t} \Big\{ & f(z_t, k_t) - i_t - \psi(i_t, k_t) 
    + \theta k_{t+1}\bigl[1 - \beta(1+R)\bigr] \\
    & + \beta\, \mathbb{E}\!\left[ V^*(z_{t+1}, k_{t+1}) \,\middle|\, z_t \right] \Big\}
\end{split}
\end{equation}
where $k_{t+1} = (1-\delta)k_t + i_t$. When $\beta(1+R) > 1$, the debt terms 
vanish ($b_{t+1} = 0$). The firm's action-value function reduces correspondingly 
to $Q^*(z_t, k_t, i_t)$, and the firm need only maintain beliefs 
$Q \sim \mathcal{GP}(\hat{Q}_t, \hat{\Sigma}_t)$ over the three-dimensional 
state-action space $(z_t, k_t, i_t)$. Proof in \ref{app:proof:debt_separation}.
\end{proposition}


\subsection{Calibration \& Implementation Details} \label{subsec:investment_calibration}

The firm's Gaussian process belief over the action-value function requires specifying 
kernel hyperparameters, raising the question of how to discipline their calibration. 
This may seem worrisome at first, as calibration requires taking a stance on the scale 
of variation in the state-action space. In this section we provide a principled basis 
for this choice.

Ultimately, the kernel encodes how quickly the firm believes the value function changes 
as it moves from $\mathbf{x}$ to $\mathbf{x}'$ in the state-action space. This belief
can be compared to the true variation in the value function, providing a measure of
how appropriately the firm generalizes from past experience to new situations.
If the firm believes the action-value function changes more slowly than it actually does,
it over-generalizes: it treats dissimilar states as interchangeable and transfers learning
too broadly. If the firm believes the action-value function changes more rapidly than it
actually does, it under-generalizes: it treats similar states as distinct and fails to
transfer learning across them. If the firm believes the action-value function changes 
at the same rate as it actually does, we consider it to be rational.
We provide a formal characterization of this (\ref{def:kernel_rationality}).

\begin{definition}[Kernel Rationality] \label{def:kernel_rationality}
An agent's GP prior $\mathcal{GP}(\hat{Q}_0, \hat{\Sigma}_0)$ is 
\textit{rational} at the steady state $\mathbf{x}^*$ if the 
believed variation in $Q^*$ equals the true variation. Define the 
rationality ratio along dimension $j$ at displacement 
$\Delta x$ as
\begin{equation}
    R_j(\Delta x) \equiv 
\frac{\operatorname{Std}\left( 
    Q^*(\mathbf{x}^* + \Delta x \, \mathbf{e}_j) 
- Q^*(\mathbf{x}^*) \mid \hat{\Sigma}_0 \right)}
    {\left| Q^*(\mathbf{x}^* + \Delta x \, \mathbf{e}_j) 
- Q^*(\mathbf{x}^*) \right|}
\end{equation}
\begin{itemize}
\item[i)] \textit{Locally rational at $\Delta x$:} 
$R_j(\Delta x) = 1$ for each dimension $j$.
\item[ii)] \textit{Globally rational:} 
$R_j(\Delta x) = 1$ for all $\Delta x > 0$ and all $j$.
\end{itemize}
When $R_j > 1$, the firm \textit{under-generalizes}: it believes $Q^*$ varies more 
than it truly does, treating similar states as distinct and failing to transfer 
learning across them. When $R_j < 1$, the firm \textit{over-generalizes}: it believes 
$Q^*$ varies less than it truly does, treating distinct states as equivalent and 
transferring learning too broadly.
\end{definition}

Note that we distinguish between local and global rationality. The true variation
in $Q^*$ is governed by the economic primitives, whereas the kernel is a 
finite-parameter object that generally cannot capture the true rate of change in $Q^*$ 
at every displacement simultaneously. For example, the action-value function for our 
firm is approximately linear in productivity $z$ and capital $k$ but quadratic in 
investment $i$ (\ref{fig:rationality_diagnostic}). The squared-exponential kernel, 
a standard choice in the GP literature and the one employed by 
\citet{IlutValchev2023}, saturates at large displacements and therefore cannot 
reproduce the true variation in $Q^*$ in each dimension $(z, k, i)$ at every 
displacement $\Delta x$. A kernel calibrated to match local variation near the steady 
state will over-generalize at large displacements, while a kernel calibrated to match large 
displacements will under-generalize locally. It is therefore prudent to calibrate the kernel at a specific 
displacement. Three choices are natural: the local neighborhood around the steady state 
(gradient matching); one standard deviation of the ergodic distribution under the 
rational policy (the typical displacement the firm experiences); and the kernel's own 
characteristic scale, where the believed standard deviation equals $\sigma_0$.

\begin{figure}[ht]
    \centering
    \includegraphics[width=1.0\textwidth]{../figures/kernel_rationality_diagnostic.pdf}
    \caption{Kernel rationality and generalization. Top row: true variation 
    $|\Delta Q^*|$ (black) versus believed variation $\hat{\sigma}_j(\Delta x)$ 
    under rational (gray), under-generalizing (red), and over-generalizing 
    (blue) calibrations. Bottom row: generalization gap 
    $G_j(\Delta x) = \hat{\sigma}_j - |\Delta Q^*|$; positive regions 
    indicate under-generalization, negative regions over-generalization. 
    Dots mark crossings where the kernel is locally rational.}
    \label{fig:rationality_diagnostic}
\end{figure}

\ref{fig:rationality_diagnostic} illustrates this limitation in the firm 
investment problem. The squared-exponential kernel (RBF) cannot simultaneously match 
the true variation in $Q^*$ at all displacements: its believed standard deviation 
is concave in $|\Delta x|$ while the true variation may be linear, concave, or 
convex depending on the dimension. At best, the believed and true variation 
curves cross at a single displacement, with the kernel under-generalizing 
($G > 0$) on one side of the crossing and over-generalizing ($G < 0$) on 
the other.
\begin{equation}
    G_j(\Delta x) \equiv \hat{\sigma}_j(\Delta x) 
    - \left| Q^*(\mathbf{x}^* + \Delta x\, \mathbf{e}_j) 
    - Q^*(\mathbf{x}^*) \right|,
\end{equation}
\ref{prop:rational_rbf} characterizes the crossing point $G_j(\Delta x) = 0$ in terms of the kernel length scale 
$\ell_j$, providing a calibration target. When the prior is exactly rational 
along dimension $j$ (so that $|\Delta Q^*_j| = \sigma_0$ at the crossing), the 
proposition yields $\ell_j = d_j / \sqrt{2 \ln 2} \approx 0.849\, d_j$, 
explaining the gap between $\ell_j$ and $d_j$ visible in the top row of 
Figure~\ref{fig:rationality_diagnostic}.

\begin{proposition}[Rational Length Scales under the RBF Kernel]
\label{prop:rational_rbf}
Under the RBF kernel $\hat{\Sigma}_0(\mathbf{x}, \mathbf{x}') = \sigma_0^2 
\exp\left(-\sum_j \frac{(x_j - x_j')^2}{2\ell_j^2}\right)$, 
the agent's GP prior is locally rational at $\mathbf{x}^*$ and 
displacement $d_j$ if and only if, for each dimension $j$, the length 
scale satisfies:
\begin{equation}
    \ell_j = \frac{d_j}{\sqrt{-2\ln\left(1 - 
    \frac{|\Delta Q^*_j|^2}{2\sigma_0^2}\right)}}
\end{equation}
where $|\Delta Q^*_j| \equiv |Q^*(\mathbf{x}^* + d_j \, \mathbf{e}_j) 
- Q^*(\mathbf{x}^*)|$ is the true variation in the action-value function 
over the chosen displacement.
\end{proposition}

\paragraph{Implementation.} We note that GP regression can be implemented via its 
kernel (dual) representation which involves inverting the $N \times N$ Gram matrix rather than 
operating in the potentially high-dimensional feature space.\footnote{This is more 
computationally efficient when the number of observations $N$ is small relative to 
the feature-space dimensionality.} In our setting the state-action space is continuous, 
and it is reasonable to maintain hundreds of observations in the firm's experience buffer. 
This representation is detailed in \ref{app:gp_regression_equivalence} and is the one used in our
implementation.